% PAGES: 119-120

We obtain that
\begin{equation}
Q \ = \ \begin{pmatrix} 0.2578 \\ 0.0075 \\ 0.2465 \\ 0.1700 \\ 0.1550 \\ 0.1580 \\ 0.0790 \end{pmatrix}.
\end{equation}
Since all the entries of $Q$ are positive, we conclude from Theorem 3.1
that the one period market model is arbitrage--free.

We conclude that the one period options market model above is both
arbitrage--free and complete.

Thus, all the other options can be priced using risk--neutral pricing;
cf. Theorem 3.4. For example, the no--arbitrage value of C1300 can be
obtained as follows:

The payoff vector $C1300_\tau$ of the call option with strike 1300 at
maturity 12/22/2012 is
\begin{equation}
C1300_\tau \ = \ \begin{pmatrix} 0 & 0 & 50 & 125 & 200 & 275 & 400 \end{pmatrix},
\end{equation}
since, e.g., $C1300_\tau(5)$ is the payoff of a call option with strike
1300 if state $\omega^5$ occurs, i.e., if $S(\tau) = 1500$, and from
(3.45) we find that
\[
C1300_\tau(5) \ = \ \max(S(\tau)-1300,0) \ = \ \max(1500-1300,0) \ = \ 200.
\]

Then, from (3.30), we find that the no--arbitrage value $C1300_{model}$
of C1300 given by this one period market model is
\begin{align*}
C1300_{model} &= C1300_\tau \, Q \ = \ \sum_{i=1}^7 Q(i) C1300_\tau(i) \\
              &= \ 139.6250,
\end{align*}
where $C1300_\tau$ is given by (3.49) and $Q$ is given by (3.48). The
corresponding percentage error is $\frac{139.6250-131.70}{131.70} =
0.0602$, i.e., $6.02\%$ over the mid value $131.70$ of C1300.

To estimate the accuracy of this one period market model, we use
formula (3.30) to find the risk--neutral value of each of the $32-7=25$
options from Table 3.1 which were not chosen as securities in the
model. These values are denoted by $V_{model}(i)$, for $i = 1:25$, and
are compared to the mid prices $V_{mid}(i)$ corresponding to the same
option, for $i = 1:25$. The percentage root--mean--squared error (RMSE)
of the model is then calculated as follows:
\begin{align*}
\text{RMSE} &= \sqrt{\frac{1}{25}\sum_{i=1}^{25} \frac{(V_{model}(i)-V_{mid}(i))^2}{V_{mid}(i)^2}} \\
            &= \ 0.0506 \ = \ 5.06\%;
\end{align*}

In other words, the options that were not chosen to be securities in
the model are priced, on average, within 5\% of their mid price by
using the seven options from this one period market model.

\section{References}

The Arrow--Debreu model was introduced in the seminal paper
\cite{5}; a modern exposition can be found in Avellaneda and Laurence
\cite{7}.

Neftci \cite{31} covers the one period market model as a general
pricing framework for derivative securities, and presents its uses by
financial engineering practitioners.

The one period binomial tree risk--neutral pricing formula (3.41) can
also be obtained by using either a hedged portfolio or payoff
replication; see Blyth \cite{8} for more details. Binomial and trinomial
trees are used in practice as numerical methods for pricing derivative
securities, and they are implemented on many time steps. The
calibration of tree models and numerical implementation details can be
found in Clewlow and Strickland \cite{11} and Wilmott \cite{45}.
